\documentclass[10pt,a4paper]{amsart}
\usepackage[margin=19mm]{geometry}
\usepackage{amsmath,amssymb,mathtools}
\usepackage[scr=rsfs,cal=euler]{mathalfa}
\usepackage{xfrac}
\usepackage[unicode,hidelinks,bookmarks=false]{hyperref}
\newcommand{\ie}{i.e.\ }
\newcommand{\cf}{cf.\ }
\newcommand{\resp}{respectively\ }
\newcommand{\Subsection}{\subsection}
\providecommand{\nobreakdash}{\nobreak\hbox{-}}
\setlength{\emergencystretch}{2em}
\multlinegap0pt
\usepackage[shortlabels]{enumitem}
\usepackage{amssymb}
\usepackage{mathtools}
\newtheorem{theorem}{Theorem}
\newtheorem{ASdescent}{Artin--Schreier descent lemma}
\newtheorem{proposition}{Proposition}
\newcommand{\Aut}{\operatorname{Aut}}
\newcommand{\Dih}{\operatorname{Dih}}
\newcommand{\Res}{\operatorname{Res}}
\title{Nakajima-extremal Artin--Schreier covers of ordinary elliptic curves in characteristic $2$}
\author{Saeed Tafazolian}
\date{Source: arXiv:2609.14830v1 (CC BY 4.0)}
\begin{document}
\maketitle

\noindent\emph{Source and licence.} Nakajima-extremal Artin--Schreier covers of ordinary elliptic curves in characteristic $2$. Version arXiv:2609.14830v1 (CC BY 4.0), distributed by arXiv under the Creative Commons Attribution 4.0 International licence, http://creativecommons.org/licenses/by/4.0/, as recorded in the arXiv licence metadata for this exact version and shown on the abstract page https://arxiv.org/abs/2609.14830v1. Full licence notice: \texttt{licenses/arxiv-2609.14830-CC-BY-4.0.txt}.\par\smallskip

\noindent\textit{Editorial scope.} Counted unit: the article's proposition prop:tower (source0.tex lines 1509--1556). The article's complete original proof of it is delivered with it (source0.tex lines 1558--1642, one \texttt{proof} environment of 85 lines). No mathematics is written, repaired or re-derived: the statement and the proof are the article's own lines, in the article's own notation, and no retained line is substituted or re-flowed.  Retained uncounted as the source-local context the proof needs: lines 227--254; lines 863--885.  Nothing was omitted from any retained span, so the package carries no omission record.  No retained line contains a citation, so the wrapper carries no bibliography and no bibliography entry.  No retained line contains a figure environment, an image-inclusion command or drawing code, so no image is delivered and the package declares no asset dependency.  Editorial formatting only, in addition: an AMS \texttt{amsart} preamble that restates the article's own declarations and macros used by the retained lines, namely the environments \texttt{theorem}, \texttt{ASdescent}, \texttt{proposition} on separate counters, together with the standard packages those lines need.  The retained environments print in this document as Theorem 1, Asdescent 1, Proposition 1 in order of appearance, whereas the article prints the same items with the numbering of its own full document.  The complete native source of the article as distributed in the arXiv e-print for this exact version is bundled unchanged as \texttt{sources/210363/source0.tex}.
\begin{theorem}\label{thm:main}
Let $q=2^h$ with $h\ge3$, let $E/k$ be an ordinary elliptic curve, let
$\omega$ be a nonzero invariant differential on $E$, and let
$\lambda\in k^\times$.  Put $T=E[q](k)$.  Choose
$f\in H^0(E,\mathcal O_E(\sum_{P\in T}[P]))$ satisfying
\[
        \Res_P(f\omega)=\lambda\qquad(P\in T),
\]
and let $X$ be the smooth projective curve with function field
$k(E)(z)$, where $z^2+z=f$.  The $k$-isomorphism class of $X$ is
independent of the choice of $f$; we denote it by
$X(E,\omega,\lambda,q)$.  This curve has the following properties:
\begin{enumerate}[label=\textup{(\roman*)}]
\item $X$ is ordinary and bielliptic, with $g(X)=q+1$;
\item the deck involution of $X\to E$ is the unique bielliptic
involution of $X$;
\item
\[
        \Aut(X)\cong \Dih(C_q)\times C_2,
        \qquad
        |\Aut(X)|=4q=4(g(X)-1);
\]
\item $\Aut(X)$ has no global fixed point on $X$, and
$(X,\Aut(X))$ is of type \textup{(ib)};
\item there is a cyclic subgroup $C\cong C_q$ acting freely on $X$, and
$X/C$ is an ordinary curve of genus $2$.
\end{enumerate}
\end{theorem}

\begin{ASdescent}\label{lem:ASdescent}
Let $k$ be an algebraically closed field of characteristic $2$, let $E/k$
be an ordinary elliptic curve, and let
\[
        T=E[q](k)\cong C_q,\qquad q=2^h,\quad h\ge1.
\]
Write $\phi:E\to F:=E/T$ for the quotient isogeny.  Let
$\pi:Y\to E$ be a connected Artin--Schreier double cover with deck
involution $u$.  Suppose that the branch locus of $\pi$ is exactly $T$,
with lower jump $1$ at every branch point, and that the translations by
$T$ lift to a subgroup $\widetilde T\leq\Aut(Y)$ such that
\[
        \widetilde T\cong T,\qquad
        \widetilde T\cap\langle u\rangle=1,
        \qquad [\widetilde T,u]=1.
\]
Then, for every nonzero invariant differential $\omega$ on $E$, there is
a unique $\lambda\in k^\times$ such that
\[
        Y\cong X(E,\omega,\lambda,q)
\]
as double covers of $E$.
\end{ASdescent}

\begin{proposition}\label{prop:tower}
Let $C=\langle\widetilde\tau\rangle\cong C_q$.  For
$0\le r\le h$, let $C_r\le C$ be the unique subgroup of order $2^r$ and
put
\[
 X_r=X/C_r.
\]
Let $K_r\le T=E[2^h](k)$ be the image of $C_r$ under the projection
$C\to T$, and put
\[
 E_r=E/K_r,\qquad q_r=2^{h-r}.
\]
Then:
\begin{enumerate}[label=\textup{(\roman*)}]
\item $C_r$ acts freely on $X$;
\item $X\to X_r$ is an unramified cyclic cover of degree $2^r$;
\item $X_r$ is ordinary and
\[
 g(X_r)=1+q_r=1+2^{h-r};
\]
\item in particular, $X/C=X_h$ is an ordinary curve of genus $2$;
\item for $0\le r\le h-1$, the quotient by the image of the deck
involution is
\[
 X_r/\langle\bar u\rangle\cong E_r,
\]
and this Artin--Schreier double cover is ramified exactly over
\[
 E_r[2^{h-r}](k),
\]
with lower jump $1$ at every branch point;
\item for $0\le r\le h-1$, choose the unique invariant differential
$\omega_r$ on $E_r$ satisfying $\phi_r^*\omega_r=\omega$, where
$\phi_r:E\to E_r$ is the quotient isogeny.  Then $X_r$ is obtained by the
same residue construction with data
\[
 (E_r,\omega_r,\lambda,q_r).
\]
In particular, if $0\le r\le h-3$, then
\[
 \Aut(X_r)\cong \Dih(C_{q_r})\times C_2,
 \qquad |\Aut(X_r)|=4q_r=4(g(X_r)-1),
\]
and $X_r$ is again of type \textup{(ib)}.
\end{enumerate}
Thus the unramified quotient tower remains inside the same non-dihedral
extremal family from genus $2^h+1$ down to genus $9$.
\end{proposition}

\begin{proof}
Every nontrivial element of $C$ projects to a nontrivial translation of
$E$, hence has no fixed point on $X$.  Thus every $C_r$ acts freely, and
$X\to X_r$ is unramified.  Riemann--Hurwitz yields
\[
 2g(X)-2=2^r\bigl(2g(X_r)-2\bigr),
\]
so
\[
 g(X_r)-1=2^{h-r}=q_r.
\]
For the $2$-rank, the Deuring--Shafarevich formula for the free
$C_r$-action gives
\[
 \gamma(X)-1=2^r\bigl(\gamma(X_r)-1\bigr).
\]
Since $\gamma(X)-1=2^h$, we obtain
\[
 \gamma(X_r)=1+2^{h-r}=g(X_r),
\]
and hence every $X_r$ is ordinary.

Because $C_r$ commutes with $u$, the involution $u$ descends to $X_r$ and
\[
 X_r/\langle\bar u\rangle
 \cong X/\langle C_r,u\rangle
 \cong E/K_r=E_r.
\]
The isogeny $\phi_r:E\to E_r$ is the quotient by a free translation
subgroup, hence is finite \emph{\'etale}.  Moreover,
$C_r\cap\langle u\rangle=1$, so $X$ is the normalization of the fibre
product
\[
 X_r\times_{E_r}E.
\]
Consequently the completed local extension for $X_r\to E_r$ becomes the
completed local extension for $X\to E$ after the unramified base change
$E\to E_r$.  Ramification groups and lower jumps are therefore unchanged.
The branch locus is the image of $T$ in $E_r$, a cyclic group of order
$q_r$.  The elliptic curve $E_r$ is ordinary, being isogenous to $E$, and
therefore this image is precisely $E_r[2^{h-r}](k)$.  Thus every branch
point has lower jump $1$.

Assume now $r\le h-1$, so $q_r$ is even.  Since $\phi_r$ is separable,
pullback on invariant differentials is an isomorphism; hence there is a
unique nonzero invariant differential $\omega_r$ on $E_r$ such that
\[
        \phi_r^*\omega_r=\omega.
\]
The cyclic group $C/C_r\cong C_{q_r}$ acts on $X_r$, commutes with the
image of the deck involution, has trivial intersection with it, and maps
isomorphically onto the translation subgroup $E_r[q_r](k)$.  Together
with the branch and jump computation above, Lemma~\ref{lem:ASdescent} applies and yields a unique $\lambda_r\in k^\times$ such that
\[
        X_r\cong X(E_r,\omega_r,\lambda_r,q_r).
\]

It remains to identify $\lambda_r$.  The curve $X$ is the normalization of
$X_r\times_{E_r}E$, and the map $\phi_r:E\to E_r$ is \emph{\'etale}.
Let $P\in E$ lie above $P_r\in E_r$.  For a local Artin--Schreier
representative $f_r$ of $X_r/E_r$, the map on completed local fields
induced by the finite \emph{\'etale} morphism $\phi_r$ has ramification
index $1$.  Hence pullback preserves residues, and
\[
 \Res_P\bigl(\phi_r^*(f_r\omega_r)\bigr)
   =\Res_{P_r}(f_r\omega_r).
\]  Since
$\phi_r^*\omega_r=\omega$ and the pullback cover is $X/E$, the left-hand
side equals $\lambda$.  Hence
\[
        \lambda_r=\lambda.
\]
Thus
\[
        X_r\cong X(E_r,\omega_r,\lambda,q_r).
\]

If $r\le h-3$, then $q_r\ge8$, so Theorem~\ref{thm:main} applies to this
new set of data.  This gives
\[
 \Aut(X_r)\cong\Dih(C_{q_r})\times C_2
\]
and shows that $X_r$ is again a non-dihedral Nakajima extremal curve of
type \textup{(ib)}.
\end{proof}

\end{document}
